\par\noindent\textbf{(1)} London关系给 \(\nabla\times\boldsymbol J=-\boldsymbol B/(\mu_0\lambda_L^2)\)，与 \(\nabla\times\boldsymbol B=\mu_0\boldsymbol J\) 联立得到 \(\nabla^2\boldsymbol B=\boldsymbol B/\lambda_L^2\)。平板一维情形 \(B_z''=B_z/\lambda_L^2\)，边界 \(B_z(\pm d/2)=B_0\)。
\par\noindent\textbf{(2)} 由两侧相同外场的偶对称性可排除 \(\sinh\) 项，解为 \(B_z(x)=B_0\cosh(x/\lambda_L)/\cosh[d/(2\lambda_L)]\)。中心磁场 \(B_z(0)=B_0\operatorname{sech}[d/(2\lambda_L)]\)，厚片 \(d\gg\lambda_L\) 时中心场指数减小。
\par\noindent\textbf{(3)} 因为 \((\nabla\times\boldsymbol B)_y=-dB_z/dx\)，得 \(J_y(x)=-(B_0/(\mu_0\lambda_L))\sinh(x/\lambda_L)/\cosh[d/(2\lambda_L)]\)。右半片积分 \(\int_0^{d/2}J_y dx=-(B_0/\mu_0)[1-\operatorname{sech}(d/(2\lambda_L))]\)，两侧电流反向。
\par\noindent\textbf{(4)} 若 \(d\ll\lambda_L\)，\(B_z(0)/B_0\simeq1-d^2/(8\lambda_L^2)\)，无法有效屏蔽平行磁场；若 \(d\gg\lambda_L\)，中心比值 \(\simeq2e^{-d/(2\lambda_L)}\)。这些结果适用于Meissner态，不包括涡旋穿透或接近临界场的非线性。
